\chapter{Foundations}

\section{Syntax and Semantics}

Consider an arbitrary programming language \(\mathcal{L}\) (a set of programs). We have:

- \textbf{Syntax}: describes which things are programs in \(\mathcal{L}\)

- \textbf{Semantics}: describes what happens when we run a program in \(\mathcal{L}\)

So, given the syntax for a programming language \(\mathcal{L}\), with a set of values \(\mathcal{V}\), we define 
\[
  \mathsf{Interpreter}_\mathcal{L} : \mathcal{L} \to \mathcal{V}
\]
Here the \textit{interpreter} takes a program in \(\mathcal{L}\) and runs it to produce a value in \(\mathcal{V}\). For example,
\[
  \mathsf{Interpreter}_\text{Python} (\verb|sorted([1, 4, 3, 2])|) = \verb|[1, 2, 3, 4]|
\]
This can be seen as one way of defining the semantics of the program.

A compiler translates from one language to another. For a program \(p\) in \(\mathcal{L}_1\) and \(\mathsf{Compiler} : \mathcal{L}_1 \to \mathcal{L}_2\) to be correct, we want 
\[
  \mathsf{Interpreter}_{\mathcal{L}_1} (p) = \mathsf{Interpreter}_{\mathcal{L}_2} (\mathsf{Compiler} (p))
\]
otherwise something is incorrect.

Notice that interpreters and compilers, in \textbf{theory}, are mathematical objects, but not in reality. They have implementations, and while we can treat implementations of programming languages as defining their semantics, it is helpful not to do this, since implementations of interpreters and compilers are themselves programs and can be buggy.

Furthermore, different implementations of the same programming language may exhibit \textit{diverging behaviors}, especially in the presence of under-specification. We can see that it is helpful to know what programs are meant to do (i.e., their semantics) separately from their implementations.